\documentclass[11pt,a4paper]{article}
% Preâmbulo comum. Não duplicar pacotes nos módulos.
\usepackage[utf8]{inputenc}
\usepackage[T1]{fontenc}
\usepackage[brazil]{babel}
\usepackage{amsmath,amssymb,amsthm,mathtools}
\usepackage{geometry}
\geometry{a4paper,margin=2.3cm}
\usepackage{microtype}
\usepackage{enumitem}
\usepackage{graphicx}
\usepackage{xcolor}
\usepackage{booktabs}
\usepackage{hyperref}
\hypersetup{colorlinks=true,linkcolor=blue!40!black,urlcolor=blue!50!black}

\setcounter{tocdepth}{2}
\numberwithin{equation}{section}

\theoremstyle{plain}
\newtheorem{teorema}{Teorema}[section]
\newtheorem{proposicao}[teorema]{Proposição}
\newtheorem{lema}[teorema]{Lema}
\newtheorem{corolario}[teorema]{Corolário}

\theoremstyle{definition}
\newtheorem{definicao}[teorema]{Definição}
\newtheorem{exemplo}[teorema]{Exemplo}
\newtheorem{exercicio}{Exercício}[section]

\theoremstyle{remark}
\newtheorem{observacao}[teorema]{Observação}
\newtheorem{contraexemplo}[teorema]{Contraexemplo}
\newtheorem{intuicao}[teorema]{Intuição}

\newcommand{\R}{\mathbb{R}}
\newcommand{\N}{\mathbb{N}}
\newcommand{\Z}{\mathbb{Z}}
\newcommand{\Q}{\mathbb{Q}}
\DeclareRobustCommand{\eps}{\varepsilon}
\DeclareRobustCommand{\vet}[1]{\vec{#1}}
\DeclareRobustCommand{\dif}{\mathrm{d}}
\DeclareRobustCommand{\abs}[1]{\left\lvert #1 \right\rvert}
\DeclareRobustCommand{\norm}[1]{\left\lVert #1 \right\rVert}
\DeclareMathOperator{\grad}{grad}
\DeclareMathOperator{\Div}{div}
\DeclareMathOperator{\rot}{rot}

\newcommand{\base}[1]{\par\medskip\noindent\textbf{Base para relembrar.} #1\par\medskip}
\newcommand{\cuidado}[1]{\par\medskip\noindent\textbf{Quando a hipótese falha.} #1\par\medskip}

\newcommand{\figura}[2]{%
  \begin{figure}[htbp]
  \centering
  \IfFileExists{#1}{%
    \includegraphics[width=0.84\textwidth]{#1}%
  }{%
    \fbox{\parbox{0.8\textwidth}{\centering Figura ainda não gerada:\\ \texttt{#1}}}%
  }
  \caption{#2}
  \end{figure}%
}

\newenvironment{esqueleto}{%
  \par\medskip\noindent\textbf{Esqueleto da demonstração.}\par\nopagebreak
  \begin{enumerate}[leftmargin=1.4em]
}{\end{enumerate}\medskip}

\title{Gabarito --- Continuidade}
\author{}
\date{}
\begin{document}
\maketitle
\tableofcontents

A ordem \'e a da se\c{c}\~ao de exerc\'icios de \texttt{02\_continuidade.tex}.

\section{B\'asico --- Linear}

\textbf{Enunciado.} $f(x)=3x-1$ \'e cont\'inua em $4$.

\begin{proof}
$f(4)=11$ e $\abs{f(x)-f(4)}=\abs{3x-1-11}=3\abs{x-4}$. Dado $\eps>0$, escolha $\delta=\eps/3$. Se $\abs{x-4}<\delta$, ent\~ao $\abs{f(x)-f(4)}<\eps$. O ponto $x=4$ entra na conta e d\'a diferen\c{c}a zero.
\end{proof}

\section{B\'asico --- Polin\^omio}

\textbf{Enunciado.} $p(x)=x^2-3x+1$ \'e cont\'inua em $2$, com $\delta$ expl\'icito.

\begin{proof}
$p(2)=4-6+1=-1$. Al\'em disso
\[
p(x)-p(2)=x^2-3x+1-(-1)=x^2-3x+2=(x-1)(x-2),
\]
logo $\abs{p(x)-p(2)}=\abs{x-2}\abs{x-1}$. Se $\abs{x-2}<1$, ent\~ao $1<x<3$, logo $0<x-1<2$ e $\abs{x-1}<2$. Portanto $\abs{p(x)-p(2)}<2\abs{x-2}$. Dado $\eps>0$, escolha $\delta=\min(1,\eps/2)$. Se $\abs{x-2}<\delta$, a diferen\c{c}a \'e menor que $2\cdot(\eps/2)=\eps$.
\end{proof}

\section{B\'asico --- Degrau}

\textbf{Enunciado.} $f(x)=0$ para $x<1$ e $f(x)=2$ para $x\ge 1$. Laterais em $1$.

\begin{proof}
Se $1-\delta<x<1$, ent\~ao $f(x)=0$, logo $\abs{f(x)-0}=0$. Qualquer $\delta>0$ mostra que o limite \`a esquerda \'e $0$. Se $1<x<1+\delta$, ent\~ao $f(x)=2$, logo o limite \`a direita \'e $2$. Como $f(1)=2$, a fun\c{c}\~ao \'e cont\'inua \`a direita em $1$ e n\~ao \'e cont\'inua \`a esquerda. Os laterais diferem, ent\~ao n\~ao h\'a continuidade bilateral.
\end{proof}

\section{B\'asico --- Extens\~ao remov\'ivel}

\textbf{Enunciado.} Estenda $(x^2-4)/(x-2)$ continuamente a $x=2$.

\begin{proof}
Para $x\neq 2$, $(x^2-4)/(x-2)=(x-2)(x+2)/(x-2)=x+2$. Ent\~ao $\abs{f(x)-4}=\abs{x-2}$. Dado $\eps>0$, $\delta=\eps$ mostra que o limite \'e $4$. Defina $f(2)=4$. O limite coincide com o valor, e a extens\~ao \'e cont\'inua em $2$.
\end{proof}

\section{Intermedi\'ario --- Sinal local}

\textbf{Enunciado.} Se $f$ \'e cont\'inua em $a$ e $f(a)>0$, ent\~ao $f$ \'e positiva numa vizinhan\c{c}a de $a$.

\begin{proof}
Seja $\eps=f(a)/2>0$. Existe $\delta>0$ tal que $\abs{x-a}<\delta$ implica $\abs{f(x)-f(a)}<\eps$. Logo
\[
f(x)>f(a)-\frac{f(a)}{2}=\frac{f(a)}{2}>0.
\]
\end{proof}

\section{Intermedi\'ario --- Raiz c\'ubica}

\textbf{Enunciado.} $x^3+x-1=0$ tem uma \'unica raiz real, e ela est\'a em $(0,1)$.

\begin{proof}
Seja $h(x)=x^3+x-1$. O polin\^omio $h$ \'e cont\'inuo em $[0,1]$. Temos $h(0)=-1<0$ e $h(1)=1>0$. Pelo teorema do valor intermedi\'ario existe $c\in(0,1)$ com $h(c)=0$.

Para a unicidade, tome $y>x$. Ent\~ao
\[
h(y)-h(x)=(y^3-x^3)+(y-x)=(y-x)(y^2+xy+x^2+1).
\]
O fator $y-x$ \'e positivo. O outro fator \'e positivo porque
\[
y^2+xy+x^2=\Bigl(x+\frac{y}{2}\Bigr)^2+\frac{3}{4}y^2\ge 0,
\]
logo $y^2+xy+x^2+1\ge 1>0$. Portanto $h(y)>h(x)$. A fun\c{c}\~ao \'e estritamente crescente em $\R$ e n\~ao pode ter duas raizes. A raiz de $(0,1)$ \'e a \'unica raiz real.
\end{proof}

\section{Intermedi\'ario --- Composi\c{c}\~ao}

\textbf{Enunciado.} $\sqrt{x^2+1}$ \'e cont\'inua em cada $a\in\R$, com $\delta'$ da raiz e $\delta$ do polin\^omio encadeados.

\begin{proof}
Seja $a\in\R$, $b=a^2+1\ge 1$ e $u(x)=x^2+1$, de modo que $f=\sqrt{\phantom{x}}\circ u$ e $u(a)=b$. Fixe $\eps>0$. Na prova da continuidade da raiz em $b>0$, a escolha
\[
\delta'=\min\Bigl(\frac{b}{2},\eps\sqrt{b}\Bigr)
\]
satisfaz: se $y\ge 0$ e $\abs{y-b}<\delta'$, ent\~ao $y>b/2$ e
\[
\abs{\sqrt{y}-\sqrt{b}}=\frac{\abs{y-b}}{\sqrt{y}+\sqrt{b}}<\frac{\abs{y-b}}{\sqrt{b}}<\eps.
\]
Resta for\c{c}ar $\abs{u(x)-b}<\delta'$. Fatore $\abs{u(x)-u(a)}=\abs{x-a}\abs{x+a}$. Se $\abs{x-a}<1$,
\[
\abs{x+a}\le \abs{x-a}+2\abs{a}<1+2\abs{a}.
\]
Seja $M=2\abs{a}+1\ge 1$. Ent\~ao $\abs{u(x)-b}<M\abs{x-a}$ sempre que $0<\abs{x-a}<1$. Escolha
\[
\delta=\min\Bigl(1,\frac{\delta'}{M}\Bigr).
\]
Se $\abs{x-a}<\delta$, ou $x=a$ e a diferen\c{c}a de $f$ \'e zero, ou $0<\abs{x-a}<\delta$ e $\abs{u(x)-b}<\delta'$. Como $u(x)\ge 1>0$, a estimativa da raiz d\'a $\abs{f(x)-f(a)}<\eps$.
\end{proof}

\section{Avan\c{c}ado --- Imagem do compacto}

\textbf{Enunciado.} Se $f\colon[a,b]\to\R$ \'e cont\'inua, a imagem \'e $[m,M]$, com $m=\min f$ e $M=\max f$.

\begin{proof}
Pelo teorema de Weierstrass existem $x_m,x_M\in[a,b]$ com $f(x_m)=m$ e $f(x_M)=M$, e $m\le f(x)\le M$ para todo $x\in[a,b]$. Logo a imagem est\'a contida em $[m,M]$, e $m$ e $M$ s\~ao valores.

Seja $k\in(m,M)$. Se $x_m<x_M$, ent\~ao $f(x_m)=m<k<M=f(x_M)$. O teorema do valor intermedi\'ario em $[x_m,x_M]$ fornece $c$ com $f(c)=k$. Se $x_M<x_m$, ent\~ao $f(x_M)>k>f(x_m)$, e o mesmo teorema, no caso em que o valor da esquerda \'e o maior, fornece $c$ em $(x_M,x_m)$. Se $x_m=x_M$, ent\~ao $m=M$ e o intervalo $(m,M)$ \'e vazio; a imagem \'e o conjunto de um ponto $[m,m]$. Em todos os casos $[m,M]$ est\'a contido na imagem.
\end{proof}

\section{Avan\c{c}ado --- M\'inimo positivo}

\textbf{Enunciado.} Cont\'inua e positiva em $[a,b]$ tem m\'inimo positivo. Em $(0,1]$, $f(x)=x$ \'e positiva e tem \'infimo $0$.

\begin{proof}
Pelo teorema de Weierstrass, o m\'inimo \'e assumido em algum $c\in[a,b]$. Por hip\'otese $f(c)>0$. O m\'inimo \'e exatamente $f(c)$, logo \'e positivo. Em particular o \'infimo, que coincide com o m\'inimo, \'e positivo.

Agora $f(x)=x$ em $(0,1]$. Para $c\in(0,1]$ e $\eps>0$, $\delta=\eps$ d\'a a continuidade: se $x\in(0,1]$ e $\abs{x-c}<\delta$, ent\~ao $\abs{f(x)-f(c)}<\eps$. Al\'em disso $f(x)>0$ em todo o dom\'inio. O n\'umero $0$ \'e cota inferior. Se $\eta>0$, o ponto $x=\min(\eta/2,1)$ est\'a em $(0,1]$ e $f(x)\le \eta/2<\eta$, logo $\eta$ n\~ao \'e cota inferior. O \'infimo \'e $0$, que n\~ao \'e positivo e n\~ao \'e assumido.
\end{proof}

\end{document}
